\begin{tikzpicture}[x=1cm, y=1cm, >=Stealth]

    % ==========================================
    % MACROS PER DISEGNARE I BLOCCHI E I DETTAGLI
    % ==========================================

    % Macro per i tensori 3D Args: x, y, width, height, depth, color, label_text
    \newcommand{\drawbox}[7]{
        \pgfmathsetmacro{\dx}{#5*0.4}
        \pgfmathsetmacro{\dy}{#5*0.4}
        % Right side (depth)
        \draw[fill=#6!60!black, draw=black!60, line width=0.2mm] (#1+#3,#2) -- ++(\dx,\dy) -- ++(0,#4) -- ++(-\dx,-\dy) -- cycle;
        % Top side
        \draw[fill=#6!80!black, draw=black!60, line width=0.2mm] (#1,#2+#4) -- ++(\dx,\dy) -- ++(#3,0) -- ++(-\dx,-\dy) -- cycle;
        % Front face
        \draw[fill=#6, draw=black!80, line width=0.2mm] (#1,#2) rectangle ++(#3,#4);

        % 2. Centered Text Node inside the box
        \node[text width=#3 cm, align=center, font=\scriptsize, text=black, transform shape] at (#1+#3/2, #2+#4/2) {#7};
    }

    \newcommand{\slicebox}[7]{
        % #1=x, #2=y, #3=w, #4=h, #5=slice_d, #6=shift_d, #7=color
        \pgfmathsetmacro{\xshift}{#6*0.4}
        \pgfmathsetmacro{\yshift}{#6*0.4}
        \pgfmathsetmacro{\dx}{#5*0.4}
        \pgfmathsetmacro{\dy}{#5*0.4}

        \coordinate (FBL) at (#1+\xshift, #2+\yshift); % Front Bottom Left della fetta
        % Right side
        \draw[fill=#7!60!black, draw=black!80, line width=0.2mm] (#1+\xshift+#3, #2+\yshift) -- ++(\dx,\dy) -- ++(0,#4) -- ++(-\dx,-\dy) -- cycle;
        % Top side
        \draw[fill=#7!80!black, draw=black!80, line width=0.2mm]  (#1+\xshift, #2+\yshift+#4) -- ++(\dx,\dy) -- ++(#3,0) -- ++(-\dx,-\dy) -- cycle;
        % Front face
        \draw[fill=#7, draw=black!80, line width=0.2mm] (FBL) rectangle ++(#3,#4);
    }

    % 3. Macro per il tensore concatenato (Bianco + Rosso + Blu)
    \newcommand{\drawsplitbox}[6]{
        % #1=x, #2=y, #3=w, #4=h, #5=d, #6=label    
        % Retro: Ultimi 12 canali (Blu)
        \slicebox{#1}{#2}{#3}{#4}{0.2*#5}{0.9*#5}{myblue}
        % Centro: Successivi 256 canali (Rossi)
        \slicebox{#1}{#2}{#3}{#4}{0.4*#5}{0.45*#5}{myred}
        % Fronte: Primi 256 canali (Bianchi)
        \slicebox{#1}{#2}{#3}{#4}{0.4*#5}{0}{mywhite}
        % Testo sulla faccia frontale
        \node[text width=#3 cm, align=center, font=\scriptsize, text=black, transform shape] at (#1+#3/2, #2+#4/2) {#6};
    }

    % 4. Macro per il Residual Layer
    \newcommand{\resblock}[3]{
    % #1=startX, #2=startY, #3=endX
    \pgfmathsetmacro{\midX}{(#1+#3)/2}
    \pgfmathsetmacro{\plusY}{#2-0.4}
    \pgfmathsetmacro{\fY}{#2-0.8}

    % Box F(x) - Componenti verdi, testo nero
    \node[draw=mygreenLine, fill=mygreen!10, thick, text=black, minimum width=0.6cm, minimum height=0.4cm, font=\scriptsize, align=center] (fx) at (\midX, \fY) {$F(x)$};
    % Nodo Sommatore (+)
    \node[draw=mygreenLine, circle, inner sep=0pt, minimum size=0.25cm, thick, text=mygreenLine, font=\scriptsize] (plus) at (#3, \plusY) {$+$};

    % Rami con frecce più piccole e raffinate (scale=0.7)
    \draw[-{Stealth[scale=0.6]}, draw=mygreenLine, thick] (#1, #2) |- (fx.west);
    \draw[-{Stealth[scale=0.6]}, draw=mygreenLine, thick] (#1, #2) |- (plus.west);
    \draw[-{Stealth[scale=0.6]}, draw=mygreenLine, thick] (fx.east) -| (plus.south);
    \draw[-{Stealth[scale=0.6]}, draw=mygreenLine, thick] (plus.north) -- (#3, #2);
    }

    % Macro per il "piano roll" (i quadratini neri)
    \newcommand{\pianoroll}[2]{
        % Background box to contain the notes
        \draw[fill=white, draw=black, line width=0.2mm] (#1, #2) rectangle ++(1.5, 1.5);
        % Note blocks
        \fill[black] (#1+0.1, #2+1.0) rectangle ++(0.2, 0.1);
        \fill[black] (#1+0.3, #2+1.1) rectangle ++(0.2, 0.1);
        \fill[black] (#1+0.5, #2+0.8) rectangle ++(0.2, 0.1);
        \fill[black] (#1+0.7, #2+0.7) rectangle ++(0.2, 0.1);
        \fill[black] (#1+0.9, #2+0.9) rectangle ++(0.2, 0.1);
        \fill[black] (#1+1.1, #2+0.6) rectangle ++(0.2, 0.1);
        \fill[black] (#1+1.3, #2+0.5) rectangle ++(0.1, 0.1);
    }

    % Colori personalizzati per i blocchi
    \colorlet{myblue}{cyan!40}
    \colorlet{myred}{red!50}
    \colorlet{mygreen}{green!50}
    \colorlet{mywhite}{white}

    % Colori personalizzati (più scuri) per le linee e le frecce
    \colorlet{myblueLine}{cyan!70!black}
    \colorlet{myredLine}{red!70!black}
    \colorlet{mygreenLine}{green!70!black}

    % ==========================================
    % PARTE SINISTRA: CONDITIONER E GENERATOR
    % ==========================================

    % --- Titoli ---
    \node at (6, 9) {\Large \textbf{Generator with conditioner}};

    % ----------- MAIN PIPELINE -------------
    % Noise z (ruotato)
    \begin{scope}[rotate around={-6:(0,1)}]
        \drawbox{0}{5.7}{1.5}{0.3}{0}{mywhite}{\textbf{Noise $z$}};
    \end{scope}
    \draw[-{Stealth}, thick, gray] (1.5, 6.2) to[out=40, in=140] node[midway, above, text=black] {\scriptsize \textsf{project and reshape}} (3.0, 6.2);

    % First tensor: 524x1x2 (Utilizziamo la nuova \drawsplitbox)
    \drawsplitbox{3}{5.5}{0.5}{0.3}{1}{}
    \resblock{3.5}{5.3}{4.85}
    \draw[-{Stealth[scale=0.8]}, thick, gray] (4, 6.3) to[out=40, in=140] node[midway, above, text=black] {\scriptsize \textsf{transp conv}} (4.9, 6.3);

    % Second tensor: 524x1x4
    \drawsplitbox{4.5}{5.5}{1}{0.3}{1}{}
    \resblock{5.5}{5.3}{6.85}
    \draw[-{Stealth[scale=0.8]}, thick, gray] (6, 6.3) to[out=40, in=140] node[midway, above, text=black] {\scriptsize \textsf{transp conv}} (7, 6.3);

    % Third tensor: 524x1x8
    \drawsplitbox{6.5}{5.5}{1.5}{0.3}{1}{}
    \resblock{8}{5.3}{9.5}
    \draw[-{Stealth[scale=0.8]}, thick, gray] (8.5, 6.3) to[out=35, in=145] node[midway, above, text=black] {\scriptsize \textsf{transp conv}} (9.5, 6.3);

    % Fourth tensor: 524x1x16
    \drawsplitbox{9}{5.5}{2}{0.3}{1}{}
    \draw[-{Stealth[scale=0.8]}, thick, gray] (11, 6.3) to[out=60, in=160] node[midway, above, text=black] {\scriptsize \textsf{transp conv}} (11.7, 6.5);

    % Final Piano Roll
    \pianoroll{11.7}{5}
    \node at (12.45, 7) {\scriptsize Output};

    % ----------- CONDITION 1D -------------    
    \begin{scope}[rotate around={-6:(0,1)}]
        \drawbox{0}{8}{1}{0.3}{0}{myblue}{\textbf{Chord}};
    \end{scope}

    \draw[-{Stealth[scale=0.8]}, thick, myblueLine] (2, 8) -| (3.7, 6.3);
    \draw[-{Stealth[scale=0.8]}, thick, myblueLine] (2, 8) -| (5.5, 6.3);
    \draw[-{Stealth[scale=0.8]}, thick, myblueLine] (2, 8) -| (7.7, 6.3);
    \draw[-{Stealth[scale=0.8]}, thick, myblueLine] (2, 8) -| (10.2, 6.3);

    \node at (6, 8.2) {\scriptsize Extend};

    % ----------- CONDITION 2D -------------
    \pianoroll{0.5}{1.8}
    \draw[-{Stealth[scale=0.8]}, thick, gray] (2.1, 1.7) to[out=-35, in=-130] node[midway, below, text=black] {\scriptsize \textsf{conv}} (3, 2.2);

    \drawbox{3}{2.3}{2}{0.3}{0.4}{myred}{}
    \draw[-{Stealth[scale=0.8]}, thick, gray] (5, 2.2) to[out=-35, in=-145] node[midway, below, text=black] {\scriptsize \textsf{conv}} (6, 2.2);
    \drawbox{6}{2.3}{1.5}{0.3}{0.4}{myred}{}
    \draw[-{Stealth[scale=0.8]}, thick, gray] (7.5, 2.2) to[out=-35, in=-145] node[midway, below, text=black] {\scriptsize \textsf{conv}} (8.5, 2.2);
    \drawbox{8.5}{2.3}{1}{0.3}{0.4}{myred}{}
    \draw[-{Stealth[scale=0.8]}, thick, gray] (9.5, 2.2) to[out=-35, in=-145] node[midway, below, text=black] {\scriptsize \textsf{conv}} (10.5, 2.2);
    \drawbox{10.5}{2.3}{0.5}{0.3}{0.4}{myred}{}

    % |- arrows
    \draw[thick, myredLine] (10.75, 2.2) |- (11.40, 1.85);
    \draw[thick, myredLine] (9, 2.2) |- (11.55, 1.7);
    \draw[thick, myredLine] (6.75, 2.2) |- (11.7, 1.55);
    \draw[thick, myredLine] (4, 2.2) |- (11.85, 1.40);

    % |- arrows
    \draw[thick, myredLine] (11.40, 1.85) |- (3.25, 3.15);
    \draw[thick, myredLine] (11.55, 1.7) |- (5.15, 3.30) ;
    \draw[thick, myredLine] (11.7, 1.55) |- (7.15, 3.45);
    \draw[thick, myredLine] (11.85, 1.40) |- (10.15, 3.60);

    % final arrows
    \draw[-{Stealth[scale=0.8]}, thick, myredLine] (3.25, 3.15) -- (3.25, 5.3);
    \draw[-{Stealth[scale=0.8]}, thick, myredLine] (5.15, 3.30) -- (5.15, 5.3);
    \draw[-{Stealth[scale=0.8]}, thick, myredLine] (7.15, 3.45) -- (7.15, 5.3);
    \draw[-{Stealth[scale=0.8]}, thick, myredLine] (10.15, 3.60) -- (10.15, 5.3);

    % ==========================================
    % SEPARATORE CENTRALE
    % ==========================================
    \draw[line width=0.4mm, gray!40] (13.5, 0.5) -- (13.5, 8);

    % ==========================================
    % PARTE DESTRA: DISCRIMINATOR
    % ==========================================

    \node at (18.5, 9) {\Large \textbf{Critic}};

    % 1. Input Piano Roll
    \pianoroll{14.5}{5}

    % 2. Converging dashed lines for Conv Layers
    \draw[dashed, thick, black!70] (16.0, 6.5) -- (17.5, 5.8);
    \draw[dashed, thick, black!70] (16.0, 5.0) -- (17.5, 5.5);
    \node at (17.25, 6.7) {\scriptsize 4 conv steps};

    % 3. Conv Feature Map Block (256x1x1)
    \drawbox{17.5}{5.5}{0.3}{0.3}{0.6}{mywhite}{}

    % 4. Minibatch Standard Deviation (sigma)
    \draw[-{Stealth[scale=0.8]}, thick, gray] (17.65, 5.5) -- (17.65, 3.7) node[below, text=black] {\scriptsize $\sigma$};
    \node[anchor=south] at (16.15, 3.35) {\scriptsize std over minibatch};

    % 5. Chord Vector Condition
    \begin{scope}[shift={(14, 0)}] % Trasla il blocco in orizzontale, allineandolo al Critic
        \begin{scope}[rotate around={-6:(0,1)}] % Usa lo stesso identico asse di rotazione del lato sinistro
            \drawbox{0}{8}{1}{0.3}{0.1}{myblue}{\textbf{Chord}};
        \end{scope}
    \end{scope}

    % 6. Reshape and Concat Arrows
    % \draw[-{Stealth[scale=0.8]}, thick, gray] (18.2, 5.75) -- (20, 5.75) node[midway, below, text=black] {\scriptsize reshape \\ \& \\ concat};
    \draw[-{Stealth[scale=0.8]}, thick, gray] (18.2, 5.75) -- (20, 5.75) node[midway, below, text=black, align=center] {\scriptsize reshape \\ \scriptsize\& \\ \scriptsize concat};
    \draw[-{Stealth[scale=0.8]}, thick, myblueLine] (16, 8) -| (20.15, 7.1);
    \draw[-{Stealth[scale=0.8]}, thick, gray] (18, 3.5) -| (20.15, 3.9);

    % 7. Combined Flattened Vector (269 dims = 256 + 1 + 12)
    \drawbox{20}{4.0}{0.3}{3.0}{0}{mywhite}{}

    % 8. Fully Connected Convergence
    \draw[dashed, thick, black!70] (20.3, 7.0) -- (22.3, 5.7);
    \draw[dashed, thick, black!70] (20.3, 4.0) -- (22.3, 5.3);
    \node at (22.15, 6.5) {\scriptsize \textsf{2 linear steps}};

    % 9. Output Validity Node
    \drawbox{22.3}{5.3}{0.4}{0.4}{0}{mywhite}{V};
    \node at (22.5, 4.8) {\scriptsize Validity};

\end{tikzpicture}